\begin{tikzpicture}
	
	%%% effect of transcription factor binding sites %%%
	\coordinate (TF effects) at (0, 0);
	
	\begin{pggfplot}[%
		at = {(TF effects)},
		yshift = -.5\baselineskip,
		total width = \textwidth,
		height = 2cm,
		xlabel = transcription factor family (motif ID),
		ylabel = condition,
		ylabel add to width = 9.164pt,
		tight y,
		enlargelimits = 0,
		grid = none,
		xticklabel style = {anchor = east, rotate = 90, inner ysep = 0pt, inner xsep = .3333em},
		xticklabels are csnames,
		facet 1/.append style = {
			pggf colorbar = {%
				horizontal,
				title = {$\log_2$(effect size)},
				title style = {yshift = .1em},
				height = .125 * \pgfkeysvalueof{/pgfplots/parent axis width},
				anchor = south west,
				at = {(parent axis.outer south west)},
				shift = {(-3\baselineskip, 2\baselineskip)},
				ytick = {-.3, 0, ..., 1},
				name = colorbar
			}
		},
		colormap 3 colors,
		y dendrogram height = .5cm,
		show y dendrogram
	]{TF_effects}
	
		\pggf_heatmap()
	
	\end{pggfplot}
	
	% mark validation TFs
	\pgfplotstableread[col sep = tab]{rawData/TF_position.tsv}\TFpostable
	
	\distance{pggfplot c1r1.west}{pggfplot c1r1.east}
	\pgfmathsetlength{\templength}{\xdistance/72}
	
	\pgfplotstableforeachcolumnelement{pos}\of\TFpostable\as\TFpos{%
		\node[anchor = south, inner xsep = 0pt] at ($(pggfplot c1r1.north west) + (-.5\templength, 0) + (\TFpos\templength, 0)$) {\textasteriskcentered};
	}
	
	
	%%% scheme of the TFBS shuffling experiment %%%
	\coordinate[yshift = -\columnsep] (shuffle scheme) at (TF effects |- xlabel.south);
	
	\node[anchor = north, font = \figsmall, text width = \fourcolumnwidth - 0.6666em, align = flush center, shift = {(.5\fourcolumnwidth, -.65)}] (seqs text) at (shuffle scheme) {50 genomic sequences with single strong binding site for a given transcription factor};

	\node[anchor = north, fill = DodgerBlue1, text = white, node font = \figsmall\bfseries, inner sep = .15em, shift = {(-.75, -.1)}] (TFBS 1) at (shuffle scheme -| seqs text) {TFBS};
	\node[anchor = south, fill = DodgerBlue2, text = white, node font = \figsmall\bfseries, inner sep = .15em, shift = {(.75, 0)}] (TFBS 3) at (seqs text.north) {TFBS};
	\node[fill = DodgerBlue3, text = white, node font = \figsmall\bfseries, inner sep = .15em] (TFBS 2) at ($(TFBS 1)!.5!(TFBS 3)$) {TFBS};
	
	\begin{pgfonlayer}{background}
		\draw[thick] (shuffle scheme |- TFBS 1) ++(.5, 0) -- ++(\fourcolumnwidth - 2cm, 0);
		\draw[thick] (shuffle scheme |- TFBS 2) ++(1, 0) -- ++(\fourcolumnwidth - 2cm, 0);
		\draw[thick] (shuffle scheme |- TFBS 3) ++(1.5, 0) -- ++(\fourcolumnwidth - 2cm, 0);
	\end{pgfonlayer}
	
	\draw[very thick, -Latex] (seqs text.south) -- ++(0, -.45) coordinate (arrow end);
	
	\node[anchor = north, font = \figsmall, text width = \fourcolumnwidth - 0.6666em, align = flush center, yshift = -.65cm] (shuffle text) at (arrow end) {replace original transcription factor binding site with randomly shuffled version};
	
	\begin{scope}
		\selectcolormodel{gray}
		\node[anchor = north, fill = DodgerBlue1, text = white, node font = \figsmall\bfseries, inner sep = .15em, shift = {(-.75, -.1)}] (TFBS 1) at (arrow end -| seqs text) {FSBT};
		\node[anchor = south, fill = DodgerBlue2, text = white, node font = \figsmall\bfseries, inner sep = .15em, shift = {(.75, 0)}] (TFBS 3) at (shuffle text.north) {TSFB};
		\node[fill = DodgerBlue3, text = white, node font = \figsmall\bfseries, inner sep = .15em] (TFBS 2) at ($(TFBS 1)!.5!(TFBS 3)$) {FBTS};
	\end{scope}
	
	\begin{pgfonlayer}{background}
		\draw[thick] (shuffle scheme |- TFBS 1) ++(.5, 0) -- ++(\fourcolumnwidth - 2cm, 0);
		\draw[thick] (shuffle scheme |- TFBS 2) ++(1, 0) -- ++(\fourcolumnwidth - 2cm, 0);
		\draw[thick] (shuffle scheme |- TFBS 3) ++(1.5, 0) -- ++(\fourcolumnwidth - 2cm, 0);
	\end{pgfonlayer}
	
	\draw[very thick, -Latex] (shuffle text.south) -- ++(0, -.45) coordinate (arrow end);
	
	\node[anchor = north, font = \figsmall, text width = \fourcolumnwidth - 0.6666em, align = flush center, yshift = -.65cm] (measure text) at (arrow end) {measure enhancer strength of sequences with WT or shuffled binding site};
	
	\node[node font = \figsmall\bfseries, align = flush center, yshift = .325cm] (Plant STARR-seq) at (measure text.north) {Plant\\STARR-seq};
	
	\coordinate[xshift = -.3cm] (leaf center) at (Plant STARR-seq.west);
	\planticon*[width = .6cm]{tobaccoLeaf} at (leaf center);
	\begin{scope}[xscale = -1]
		\planticon[width = .323cm]{pipe} at ($(leaf center) + (.35, .16)$);
	\end{scope}
	
	\coordinate[xshift = .3cm] (protoplasts) at (Plant STARR-seq.east);
	\planticon[width = .3cm]{protoplast} at ($(protoplasts) + (90:.15)$);
	\planticon[width = .3cm]{protoplast} at ($(protoplasts) + (162:.15)$);
	\planticon[width = .3cm]{protoplast} at ($(protoplasts) + (234:.15)$);
	\planticon[width = .3cm]{protoplast} at ($(protoplasts) + (306:.15)$);
	\planticon[width = .3cm]{protoplast} at ($(protoplasts) + (18:.15)$);
	\planticon[width = .3cm, rotate = 15]{maize} at ($(protoplasts) + (.375, -.1)$);
	
	
	%%% comparison of TFBS shuffle effect with main library %%%
	\coordinate (correlation) at (shuffle scheme -| \textwidth - \threequartercolumnwidth, 0);
	
	\begin{pggfplot}[%
		at = {(correlation)},
		total width = \threequartercolumnwidth,
		xlabel = {\textbf{transcription factor binding site shuffling:} $\log_2$(effect size)},
		ylabel = {\textbf{ACR library:} $\log_2$(effect size)},
		title is csname,
		title color from name,
		xytick = {-1.2, -.9, ..., 1.2}
	]{TFBS_shuffle_tamsACR}
	
		\pggf_annotate(diagonal)
	
		\pggf_trendline()
	
		\pggf_scatter(
			color from column name,
			mark size = 1.5
		)
	
		\pggf_stats(
			stats = {n,spearman,rsquare}
		)
	
	\end{pggfplot}
	
	
	%%% scheme of synthetic enhancer design %%%
	\coordinate[yshift = -\columnsep] (synTF scheme) at (TF effects |- xlabel.south);
	
	\pgfplotsset{letter colors = DNA}
	
	\pgfmathsetlength{\templength}{.45cm}
	
	\coordinate[yshift = -\columnsep - 2\templength] (TFs) at (synTF scheme);
	
	% nucleotide composition
	\coordinate[shift = {(.5\columnsep + \templength, 1.5\templength)}] (low GC) at (TFs);
	\coordinate[shift = {(.5\columnsep + \templength, -1.5\templength)}] (high GC) at (TFs);
	
	\pgfplotstableread[col sep = tab]{rawData/synEnh_nucFreq.tsv}\nucFreqTable
	\pgfplotstableforeachcolumn\nucFreqTable\as\GCbin{
		\pgfplotstablegetelem{0}{\GCbin}\of\nucFreqTable
		\foreach [expand list, evaluate = \y as \x using 180 - \y, remember = \x as \lastx (initially 180)] \base/\y in {\pgfplotsretval}{
			\fill[letter\base col!50] (\GCbin) -- ++(\lastx:\templength) arc[start angle = \lastx, end angle = \x, radius = \templength] coordinate[pos = .5] (c\base) -- cycle;
			\node[font = \figsmall] at ($(\GCbin)!.6!(c\base)$) {\base};
		}
	}
	
	% sequences
	\coordinate[xshift = .5\columnsep + \templength] (seqs) at (TFs -| low GC);
	\coordinate (low GC seq) at (low GC -| seqs);
	\coordinate (high GC seq) at (high GC -| seqs);
	
	\distance{seqs}{\threecolumnwidth - .5\columnsep, 0}
	
	\foreach \x/\y in {low GC/south, high GC/north}{
		\draw[line width = .06cm, \x] (\x\space seq) node[anchor = \y\space west, text = black, text depth = 0pt, xshift = -.5\columnsep] {25 sequences with \textbf{\x} content} -- ++(\xdistance, 0) coordinate[pos = .165] (\x\space 1) coordinate[pos = .5] (\x\space 2) coordinate[pos = .835] (\x\space 3) coordinate (\x\space seq end);
	}
	
	% TFs
	\foreach \x in {1, 2, 3}{
		\node[fill = DodgerBlue1, text = white, node font = \figsmall\bfseries, inner sep = .15em] (TFBS \x) at (TFs -| low GC \x) {TFBS};
		
		\draw[-Latex, shorten > = .08cm, DodgerBlue1!50] (TFBS \x.north) ++(0, .05) -- (low GC \x);
		\draw[-Latex, shorten > = .08cm, DodgerBlue1!50] (TFBS \x.south) ++(0, -.05) -- (high GC \x);
	}
	
	\node[anchor = base, node font = \figsmall] at ($(TFBS 1.base)!.5!(TFBS 2.base)$) {and/or};
	\node[anchor = base, node font = \figsmall] at ($(TFBS 2.base)!.5!(TFBS 3.base)$) {and/or};


	% tested TFBSs
	\coordinate[yshift = -.5\columnsep - \templength] (TFBSs) at (synTF scheme |- high GC);
	
	\node[anchor = north, font = \bfseries] (title) at ($(TFBSs)!.5!(TFBSs -| low GC seq end)$) {transcription factor binding sites (TFBS)\settowidth{\global\templength}{\figsmall\usename{TF_13}}};

	\pgfplotstableread[col sep = tab]{rawData/synEnh_TFs.tsv}\synEnhTFtable

	\pgfplotstableforeachcolumnelement{col1}\of\synEnhTFtable\as\x{
		\pgfmathsetmacro{\xi}{int(\pgfplotstablerow + 1)}
		\node[anchor = base west, node font = \figsmall, yshift = -\xi*1.15\baselineskip, text width = \templength, align = right] at (TFBSs |- title.base) {\usename{\x}};
	}
	\pgfplotstableforeachcolumnelement{col2}\of\synEnhTFtable\as\x{
		\pgfmathsetmacro{\xi}{int(\pgfplotstablerow + 1)}
		\node[anchor = base east, node font = \figsmall, yshift = -\xi*1.15\baselineskip] (TFBS name \xi) at (low GC seq end |- title.base) {\usename{\x}};
	}

	% connect TFBS list with TFBSs
	\draw[decorate, decoration = {brace, amplitude = .3cm}] ($(TFBS name 1.north east) + (-.1, 0)$) -- ($(TFBS name 6.south east) + (-.1, 0)$) coordinate[pos = .5, xshift = .3cm - \pgflinewidth] (curly);
	\draw[thick, -Latex, shorten > = .05cm, rounded corners = .15cm] (curly) -- ++(.2, 0) |- (TFBS 3.east);


	%%% strength of synthetic enhancers by number of inserted TFBSs %%%
	\coordinate (synEnh strength) at (synTF scheme -| \textwidth - \twothirdcolumnwidth, 0);
	
	% style to shift cld labels
	\tikzset{
		label shift/.code = {
			\ifnum1=#1\relax
				\pgfkeysalso{xshift = .2em}
			\fi
			\def\tmpA{maize}
			\ifx\tmpA\pggfcolname\relax
				\ifnum1=#1\relax
					\pgfkeysalso{yshift = .5ex}
				\fi
				\ifnum2=#1\relax
					\pgfkeysalso{yshift = 1.5ex}
				\fi
			\fi
		}
	}
	
	% group styles
	\pgfplotsset{
		low GC/.code = {
			\def\tmpA{dark}
			\ifx\tmpA\pggfcolname\relax
				\pgfkeysalso{fill = \pggfcolname!50!white}
			\else
				\pgfkeysalso{fill = \pggfcolname}
			\fi
		},
		high GC/.code = {
			\def\tmpA{dark}
			\ifx\tmpA\pggfcolname\relax
				\pgfkeysalso{fill = \pggfcolname}
			\else
				\pgfkeysalso{fill = \pggfcolname!50!black}
			\fi
		}
	}
	
	\begin{pggfplot}[%
		at = {(synEnh strength)},
		total width = \twothirdcolumnwidth,
		xlabel = number of inserted transcription factor binding sites,
		ylabel = \enrichment,
		enlarge x limits = 0,
		enlarge y limits = {lower = 0.25, upper = 0.15},
		title is csname,
		title color from name,
		legend pos = north west,
		ytick = {-12, -10, ..., 12},
		legend columns = 3,
		legend style = {anchor = north, at = {(.5, 1)}},
		legend image width = .75\baselineskip
	]{synEnh_strength}
		
		\pggf_annotate(
			manual legend = {
				{\hspace*{-\groupplotsep}\textbf{GC:}\hspace*{-\groupplotsep}} = {empty legend},
				low = {boxplot legend = y, fill = gray!25!white},
				high = {boxplot legend = y, fill = gray!75!black}
			},
			only facet = 1
		)
		
		\pggf_hline(densely dotted)
	
		\pggf_boxplot(
			style from column = {GC_bin},
			sample size = {anchor = west, rotate = 90},
			cld = {anchor = south, text depth = 0pt, label shift = #1}{max},
			forget plot
		)
	
	\end{pggfplot}
	
	
	%%% subfigure labels
	\subfiglabel[yshift = .5\baselineskip]{TF effects}
	\subfiglabel{shuffle scheme}
	\subfiglabel{correlation}
	\subfiglabel{synTF scheme}
	\subfiglabel{synEnh strength}

\end{tikzpicture}